\documentclass[11pt]{article}

\usepackage[T1]{fontenc}
\usepackage[margin=1in]{geometry}
\usepackage{amsmath,amssymb,amsthm}
\usepackage{stmaryrd}
\usepackage{array}
\usepackage{booktabs}
\usepackage[colorlinks=true,linkcolor=blue,citecolor=blue,urlcolor=blue]{hyperref}

% ---- theorem environments -------------------------------------------------
\theoremstyle{plain}
\newtheorem{theorem}{Theorem}
\newtheorem{proposition}[theorem]{Proposition}
\newtheorem{lemma}[theorem]{Lemma}
\newtheorem{conjecture}[theorem]{Conjecture}
\theoremstyle{definition}
\newtheorem{definition}[theorem]{Definition}
\newtheorem{question}[theorem]{Question}
\theoremstyle{remark}
\newtheorem{remark}[theorem]{Remark}

% ---- notation -------------------------------------------------------------
\newcommand{\Set}{\mathsf{Set}}
\newcommand{\Cat}{\mathsf{Cat}}
\newcommand{\Cont}{\mathsf{Cont}}
\newcommand{\DCont}{\mathsf{DCont}}
\newcommand{\Vecfd}{\mathsf{Vec}_{\mathrm{fd}}}
\newcommand{\OrgTr}{\mathsf{OrgTr}}
\newcommand{\Sk}{\mathrm{Sk}}
\newcommand{\Aut}{\mathrm{Aut}}
\newcommand{\op}{\mathrm{op}}
\newcommand{\id}{\mathrm{id}}
\newcommand{\comp}{\mathbin{\triangleleft}}
\newcommand{\grade}[1]{\textsf{\small\textbf{#1}}}

% ---- front matter ---------------------------------------------------------
\title{Total composition \textsc{constructs}; partial composition
       \textsc{lifts}\\[2pt]
       \large A speculative mechanism, a failed degree prediction,
       and a name for the failure}
\author{MacBeth \\ \small Kodamai}
\date{31 August 2026}

\begin{document}
\maketitle

\begin{center}
\fbox{\begin{minipage}{0.92\textwidth}
\small
\textbf{Recipient:} Rick.\\[2pt]
\textbf{Author:} MacBeth (Kodamai). \qquad \textbf{Date:} 2026-08-31.\\[2pt]
This note corresponds to \texttt{work-in-progress} commit
\texttt{f95140a8\ldots} of the repository \texttt{scotmacbeth/work-in-progress};
in full,
\begin{center}
\texttt{f95140a88a390dfb6c53694535c404f189500f9b}.
\end{center}
Everything asserted below is graded. \grade{proved} means I have a proof I stand
behind; \grade{computed} means a finite verification; \grade{speculative} means a
proposal with no proof. The central thesis of this note is \grade{speculative} and
its most attractive consequence is \grade{speculative and currently unconfirmed}.
Please read the grades as binding, not decorative.
\end{minipage}}
\end{center}

\begin{abstract}
Rick --- a short note on something adjacent to your front. I have a four-instance
pattern separating compositional primitives into those that are \emph{total} and
those that are \emph{partial}, and a proposed mechanism: total composition
\textsc{constructs} the composite by a universal property, partial composition
\textsc{solves a lifting problem}, and only lifting problems carry cohomology. Our
shared object --- the Zappa--Sz\'ep obstruction $[\omega]\in
H^2(\Sk_{\mathcal C};\mathcal D)$ --- is the one data point in the pattern that is
\grade{proved}. The mechanism predicts a companion obstruction in $H^1$ for
sheaf-theoretic multi-agent composition, one degree below ours. I tested that
prediction against the literature and it is \textbf{not confirmed}; the most
tempting near-miss is a genuine $H^1$ theorem on the wrong axis. I name that failure
mode a \textsc{collision} and give its diagnostic. Three questions for you at the end.
\end{abstract}

%============================================================================
\section{The thesis}\label{sec:thesis}
%============================================================================

\noindent\fbox{\begin{minipage}{0.96\textwidth}
\textbf{Thesis} (\grade{speculative}).
\emph{Composition is total exactly when the composite is \textsc{constructed} by a
universal property, and partial exactly when the composite is a \textsc{solution to
a lifting problem}. Colimits always exist, given a global cocompleteness precondition
checked once; lifts need not, and the obstruction to a lift is a cohomology class.}
\end{minipage}}

\medskip

I want to be exact about what is and is not being claimed. I have \emph{not} proved
``universal-property-constructed $\Rightarrow$ total'', nor ``lifting problem
$\Rightarrow$ cohomological obstruction'', in any generality. Both are
folklore-shaped statements that I am using as a \emph{lens}, not as lemmas. What I
have is four data points that the lens sorts correctly, one visible symptom that the
lens explains (\S\ref{sec:asym}), and one falsifiable prediction that the lens makes
(\S\ref{sec:pred}) --- which is where the interesting part begins, because the
prediction did not survive contact with the literature and the way it failed is
itself worth a name.

Your instincts here are better than mine, which is why you are getting this rather
than a finished paper.

%============================================================================
\section{Four data points}\label{sec:data}
%============================================================================

\begin{center}
\small
\begin{tabular}{@{}p{0.30\textwidth}p{0.11\textwidth}p{0.25\textwidth}p{0.24\textwidth}@{}}
\toprule
\textbf{compositional primitive} & \textbf{total?} & \textbf{obstruction} &
\textbf{status / source} \\
\midrule
\textbf{(1)} pushout / structured cospans
  & total
  & none found anywhere in the literature
  & \emph{checked absence} of obstruction language; browse 2026-09-07,
    \texttt{2509.18475}, \texttt{2301.01445} \\[4pt]
\textbf{(2)} Spivak's $\#$ on $\OrgTr$
  & total, \emph{unconditionally}
  & none
  & \texttt{2602.17917}, Prop.\ 4.3 \\[4pt]
\textbf{(3)} Zappa--Sz\'ep / matched pairs of \emph{directed containers}
  & \textbf{partial}
  & $[\omega]\in H^2(\Sk_{\mathcal C};\mathcal D)$
  & \grade{proved} (mine) \\[4pt]
\textbf{(4)} sheaf gluing for multi-agent systems
  & \emph{the test case}
  & \textbf{predicted} $H^1$
  & \grade{speculative}, \textbf{not confirmed}; see \S\ref{sec:pred} \\
\bottomrule
\end{tabular}
\end{center}

\medskip

\noindent\textbf{(1) Pushouts and structured cospans.} The Baez--Courser--Fong line,
as carried forward in the ACSet/structured-cospan tooling, composes by pushout. The
data point is an \emph{absence}, and I want it labelled as such: I went looking for
obstruction language in that literature and did not find it. That is evidence that
nobody \emph{needed} one, not a theorem that none exists. Nothing can fail per-pair
in a pushout: the composite is built out of the two given objects, and the only
precondition --- cocompleteness of the ambient category --- is discharged once, for
the whole setting, before any particular pair is considered.

\medskip

\noindent\textbf{(2) Spivak's $\#$ on organisational trees.} I had expected this one
to go the other way, and it does not. The composition operator $\#$ on $\OrgTr$ is
\emph{unconditionally total} (\texttt{2602.17917}, Prop.\ 4.3). I record this
against my own earlier hope that a class $[\omega]$ would appear one level up in the
orchestration hierarchy; it does not, and my $[\omega]$ stays a
directed-container phenomenon. Second construction-side data point.

\medskip

\noindent\textbf{(3) Zappa--Sz\'ep products of directed containers.} \grade{proved},
and your point of contact with all of this. For $K=\mathcal C\bowtie\mathcal D$ a
strict factorisation, my pairwise criterion is a conjunction $(L)\wedge(G)$ of a
freeness condition on hom-presheaves and a global closure condition; and
\[
  (G)\iff [\omega]=0 \in H^2(\Sk_{\mathcal C};\,\mathcal D).
\]
The tower underneath is citation --- Rosebrugh--Wood for existence, Baues--Wirsching
for the abelian classification, Pirashvili for the non-abelian case; my contribution
is the identification and the rigid-twist ($\mathbb Z/2$) computation. Note the
shape: this composition does not \emph{build} $K$. It asks whether a \emph{given}
$K$ factors --- a lifting problem against the skeleton $\Sk_{\mathcal C}$, and
lifting problems against a skeleton are exactly what $H^2$ measures.

The blockchain re-entrancy instance is the smallest witness I have that the class is
not vacuous: there $[\omega]=\varepsilon$ with $H^2\cong\mathbb F_2$, machine-checked
in Lean.

\medskip

\noindent\textbf{(4) Sheaf gluing.} The test case, and \S\ref{sec:pred} is about it.

%============================================================================
\section{The visible symptom: a precondition asymmetry}\label{sec:asym}
%============================================================================

Before the prediction, the reason I believe there is \emph{something} here rather
than a tautology. The two families differ in an observable way that has nothing to do
with which tensor they use, and everything to do with \emph{when their preconditions
are discharged}:

\begin{center}
\begin{tabular}{@{}ll@{}}
\toprule
\textbf{total composition} & \textbf{partial composition} \\
\midrule
precondition is \textbf{global} & precondition is \textbf{per-pair} \\
checked \textbf{once}, for the ambient category & checked \textbf{again}, for every pair \\
shape: ``is $\mathcal E$ cocomplete?'' & shape: ``does $[\omega]$ vanish for \emph{this} $(\mathcal C,\mathcal D)$?'' \\
failure is a \emph{well-formedness} error & failure is \emph{structural} \\
\bottomrule
\end{tabular}
\end{center}

\medskip

This is the part I would defend hardest, because it is empirical rather than
theoretical: the two schools of applied category theory really do differ in the
\emph{kind} of hypothesis they carry, not merely in the monoidal structure they
carry it on. The colimit school's hypotheses are ambient; the Zappa--Sz\'ep school's
are pairwise and cohomological. The thesis of \S\ref{sec:thesis} is my attempt to say
\emph{why}, and the asymmetry is the thing it is trying to explain.

There is a consequence a systems engineer can use before modelling anything: \emph{does
your composition operator build the composite, or does it look for one?} Build
$\Rightarrow$ composition is always defined and correctness is a global
well-formedness check. Look $\Rightarrow$ composition is partial, failures are
structural rather than accidental, and there is in principle a computable invariant
saying in advance which pairs compose. Supply chains, re-entrancy, and agent
orchestration all sit on the \emph{lifting} side --- which is exactly why
compositional failure there has economic consequences, and why compositional failure
in a pushout-based model does not exist to have consequences.

%============================================================================
\section{The prediction, and its honest status}\label{sec:pred}
%============================================================================

Sheaf gluing is a \emph{limit} (matching families), not a colimit, and its totality
is conditional on a \emph{cocycle condition on a cover}. So the mechanism does not
predict ``sheaves are total''. It predicts:

\medskip
\noindent\fbox{\begin{minipage}{0.96\textwidth}
\textbf{Prediction} (\grade{speculative and currently unconfirmed}).
\emph{Sheaf-theoretic multi-agent composition carries an obstruction, and it sits in
$H^1$ --- one degree below the Zappa--Sz\'ep $H^2$ --- because gluing local sections
is a descent problem (torsor-flavoured, $H^1$), whereas Zappa--Sz\'ep asks for an
extension/factorisation ($H^2$).}
\end{minipage}}
\medskip

I tested this. The literature pass is \textbf{done}, and it returned
\emph{not confirmed}. Three sources are relevant, all deep-read:

\begin{center}
\small
\begin{tabular}{@{}p{0.16\textwidth}p{0.34\textwidth}p{0.08\textwidth}p{0.34\textwidth}@{}}
\toprule
\textbf{source} & \textbf{what it does} & \textbf{degree} & \textbf{verdict against the prediction} \\
\midrule
\texttt{2606.01663} & Nash equilibria as \textbf{global sections} in a topos & $H^0$ &
wrong degree; no gluing obstruction claimed \\[4pt]
\texttt{2605.01879} & plan-incoherence \emph{asserted} as gluing failure; ``$H^1/H^2$
obstructions to a global plan'' & $H^1$/$H^2$ &
right degree, \textbf{no theorem is proved} --- appetite, not evidence \\[4pt]
\texttt{2605.11204} \newline (Anwer--Riess--Hale), \textbf{Thm 2} & edge potentials of
conservative, edge-separable forces recoverable from trajectory data iff
$H^1(G;\mathcal F)=0$ & $H^1$ &
a genuine $H^1$ theorem, \textbf{wrong axis} --- see \S\ref{sec:collision} \\
\bottomrule
\end{tabular}
\end{center}

\medskip

The third row is the dangerous one, and it is the reason this note exists. Anwer--
Riess--Hale prove a real theorem in exactly the predicted degree. Citing it as
confirmation would have been wrong. Their $H^1$ measures \textbf{identifiability}:
can the local interaction laws be recovered from observed trajectories? Mine would
measure \textbf{existence of a glued composite}: does a family of locally defined
compositions assemble into a global one? Same degree; no shared mechanism; different
site (a communication graph versus a cover of a directed container); no shared
functor of which the two $H^1$'s are shadows.

I state the verdict plainly: \textbf{the degree prediction is unconfirmed, and
nothing in the literature I have read confirms it.} What \texttt{2605.01879} shows is
only that $H^1$ is the degree practitioners reach for when they gesture at this; that
is sociology, not evidence.

%============================================================================
\section{Collision --- a name for this failure mode}\label{sec:collision}
%============================================================================

I have been keeping a small taxonomy of ways a distinction stops being visible.
\textsc{Fusion} lives in the base: two logically independent conditions coincide over
$\Set$ and come apart over $\Vecfd$. \textsc{Identification} lives in the
representation: two non-isomorphic objects present the same functor, so the question
stops referring. This section is about the third, which is the one that costs me
time, because it is invisible to exactly the summaries I write.

\begin{definition}[Collision]\label{def:collision}
A \emph{collision} occurs when a \emph{lossy invariant} --- a cardinality, a
cohomological degree, a count --- takes the same value on two structures that
genuinely differ. Nothing has fused and nothing has been identified: the objects are
perfectly distinct; the \emph{summary statistic} chosen to describe them is not
injective, and the agreement therefore carries no information.
\end{definition}

\begin{remark}[Where each mode lives]
\textsc{Fusion} lives in the base --- change the base and it goes away.
\textsc{Identification} lives in the representation --- no change of base helps; a
presentation must be fixed by fiat. \textsc{Collision} lives in \emph{my description
of the objects}. Fusion and identification are facts about mathematics; collision is
a fact about compression --- which is why it is the mode my own note-taking
manufactures, and the only one of the three I control.
\end{remark}

\begin{quote}
\textbf{Diagnostic}, to run before any ``same invariant, therefore same phenomenon''
claim: \emph{name the two structures; name the invariant; ask what functor the
invariant is a shadow of.} If the two shadows come from different functors ---
cardinality of a hom-set versus of a colimit; $H^1$ of identifiability versus $H^1$ of
descent --- the agreement is a collision and carries no information.
\end{quote}

\noindent\textbf{The sister instance, for calibration.} The reason I trust the
diagnostic is that I have the same error in a setting where the ground truth is
finite and checkable. In my counterexample to a conjectured left adjoint over
$\Vecfd$ (\texttt{proofs/2026-08-30-left-adjoint-over-vec.md}), at
$\dim P_s=1$ with $|T|=2$ over $\mathbb F_2$, the canonical comparison map $\kappa$
has \textbf{4 elements on each side} --- and is still not a bijection. It
double-counts the zero map and misses $e_0+e_1$. \grade{computed}. A cardinality-only
check passes a failing adjunction; the map had to be examined.

\begin{quote}
\textbf{Check the map, not the count; check the axis, not the degree.}
\end{quote}

\noindent Anwer--Riess--Hale is the degree-shaped version of the $4=4$ coincidence.
That is the whole content of \S\ref{sec:pred}'s negative verdict.

%============================================================================
\section{Where this stands}\label{sec:stands}
%============================================================================

The literature pass is finished and its answer is \emph{nobody has proved this}. So
the outstanding work is not another browse; it is a proof attempt on my own model.
The question to pose formally:

\begin{question}\label{q:descent}
Is multi-agent composition over a \emph{cover of a directed container} a descent
problem --- and if so, does its obstruction live in $H^1(\text{cover};\Aut)$?
\end{question}

Two honest caveats attached to Question~\ref{q:descent}. First, ``cover of a directed
container'' is not yet a definition I have written down; making it one is most of the
work, and it may well be that the natural notion of cover makes the gluing total, in
which case data point (4) joins (1) and (2) on the construction side and the
\emph{degree} half of the thesis is simply vacuous while the coarse total/partial
split survives. Second, if an obstruction does appear, I must state its axis before
comparing it to anything --- by Definition~\ref{def:collision}, a landing in degree
$1$ tells me nothing on its own.

%============================================================================
\section{What I am asking}\label{sec:asking}
%============================================================================

Three questions, in order of how much I would value an answer.

\begin{enumerate}
\item \textbf{Is the dichotomy real?} You have the obstruction-theory instincts here
and I do not. Does the total/partial split look like a genuine structural dichotomy
--- or is it a restatement of ``limits are unique, extensions are not'', dressed up
with four examples? I would rather hear that it is folklore with a different name
than publish it as a discovery.

\item \textbf{Am I committing my own collision error?} The degree prediction --- $H^1$
for gluing, $H^2$ for Zappa--Sz\'ep --- is the part I like most, which is exactly why
I distrust it. Is there a principled reason a descent obstruction should sit one
degree below a factorisation obstruction, or am I pattern-matching on degrees in
precisely the way \S\ref{sec:collision} forbids? A verdict of ``you are doing the
thing you just named'' is a completely acceptable answer and I would like it stated
bluntly.

\item \textbf{Do you know a worked case?} A composition obstruction --- existence of a
composite, not identifiability, not recovery, not classification of twists on an
already-existing object --- genuinely landing in $H^1$. If one exists anywhere in your
territory, it settles the degree question faster than my proof attempt will. If your
answer is that composition obstructions are $H^2$-shaped by their nature and $H^1$ is
the wrong place to look, that is equally decisive and I will drop the degree half of
the thesis.
\end{enumerate}

\medskip

\noindent Nothing here needs a fast reply; the note is a snapshot, not a request for
a referee report. But the second question is the one I cannot answer from inside my
own framing, and you are the person who can.

\medskip
\noindent --- MacBeth

\end{document}
